% Test of keeping listings together (didactic): in the notes (memoir +
% beamerarticle), a listing of at most keeplines lines doesn't break across
% pages, a longer one leaves at least keepends lines on each side of a
% break, and the paragraph that leads into a listing stays with it.
% Each case starts on a new page, after a rule that leaves room for a few
% lines only, so that without didactic the listing would break there.
%
% Build: see the Makefile (three runs, for the page references and case N).  The
% document checks itself: the last page lists every check with PASS or
% FAIL, and the log has a line "didactic keep test: <check> PASS/FAIL"
% for each; the Makefile fails on any FAIL.
%
% Expected in the resulting PDF, besides PASS everywhere:
%   A  a short listing (10 lines) with room for 4: the lead-in and the
%      whole listing on the next page, the first page empty below the rule.
%   B  a long listing (30 lines) with room for the lead-in and 2 lines:
%      the lead-in and the first three lines on the next page.
%   C  a long listing with room for 10 lines: it breaks there.
%   D  a long listing with room for all but one line: the last three
%      lines on the next page, not one.
%   I  as D, but line 28 is empty: lines 27 to 30 on the next page, so
%      that three lines that aren't empty are there and the page doesn't
%      start with the empty line.
%   E  a short listing read by \inputminted (5 lines), with room for 3
%      lines: on the next page with its lead-in.
%   F  a code chunk woven for minted (as noweb.sty's [minted] option and
%      the tominted weave give), with room for 2 lines: the lead-in, the
%      header line "<hello.py>=" and the code on the next page.
%      Skipped without noweb.sty.
%   J  a pair of didacticpair (as \codebycode and \runbyrun make), 8 lines
%      tall, with room for 2 lines: on the next page with its lead-in.
%   L  a short listing (10 lines) with room for 8 of its lines: moving it
%      would leave more than keepmaxgap (5) lines empty, so it breaks, with
%      at least three lines on each side.
%   K  a section heading followed by four lines of text, with room for it
%      and two lines: it moves to the next page with three lines
%      (keepafterheading).
%   M  the heading of a block (example), with room for two of its lines:
%      the block moves to the next page.
%   N  a section heading, a sentence and a longtable whose caption and
%      first row don't fit: the heading moves to the page of the table, in
%      the run after the one that finds it (so three runs).
%   G  a short listing in an item of a list: the item's label "1." beside
%      its first line, as without didactic.
%   H  the vertical space around a listing is as without didactic: the
%      text after it starts at the same height as without (not checked
%      automatically; compare with keeplines=0, keepends=1,
%      keepleadin=false).
\documentclass[a4paper]{memoir}
\usepackage[british]{babel}
\usepackage{amsmath,amssymb,amsthm}
\usepackage{didactic}
\usepackage[noamsthm,notheorems]{beamerarticle}
\usepackage{refcount}
\usepackage{longtable}
\IfFileExists{noweb.sty}{\usepackage[minted]{noweb}}{}

\makeatletter
% \roomfor{n}: start a new page and fill it so that n lines of text fit
% under the rule.
\newcommand\roomfor[1]{%
  \clearpage
  \noindent\rule{\linewidth}{0.4pt}\par
  \vspace*{\dimexpr\textheight-\topskip-#1\baselineskip-\baselineskip\relax}%
  \par}
% \testfile{prefix}{n}{before}{after}: write the file \jobname-prefix.txt
% of n lines "line <i>", each followed by
% <before>\label{prefix L<i>}<after>, except an empty line set by
% \testblank.
\newcount\@keeptest@i
\newcommand\testfile[4]{%
  \immediate\openout\@keeptest@out=\jobname-#1.txt
  \@keeptest@i\z@
  \loop
    \advance\@keeptest@i\@ne
    \ifnum\@keeptest@i=\@keeptest@blank
      \immediate\write\@keeptest@out{}%
    \else
      \immediate\write\@keeptest@out{%
        line \the\@keeptest@i#3\string\label{#1L\the\@keeptest@i}#4}%
    \fi
  \ifnum\@keeptest@i<#2\repeat
  \immediate\closeout\@keeptest@out
  \global\@keeptest@blank\z@}
% \testblank{i}: line i of the next file is empty.
\newcount\@keeptest@blank
\newcommand\testblank[1]{\global\@keeptest@blank=#1\relax}
% \testlisting{prefix}{n}: such a file as a listing of fancyvrb.
\newcommand\testlisting[2]{%
  \testfile{#1}{#2}{}{}%
  \VerbatimInput[commandchars=\\\{\}]{\jobname-#1.txt}}
\newwrite\@keeptest@out
% \testcheck{name}{label1}{=|<}{label2}: are the labels on the same page (=)
% or is the first on an earlier page (<)?
\newcommand\testcheck[4]{%
  \ifnum\getpagerefnumber{#2}#3\getpagerefnumber{#4}\relax
    \def\@keeptest@result{PASS}%
  \else
    \def\@keeptest@result{FAIL}%
  \fi
  \typeout{didactic keep test: #1 \@keeptest@result}%
  \item #1: \@keeptest@result\ (#2 on page \getpagerefnumber{#2}, #4 on
    page \getpagerefnumber{#4})}
\makeatother

\begin{document}

\roomfor{4}
\section*{A: a short listing}
Lead-in A, which says what the listing shows.\label{a-lead}
\testlisting{a}{10}
Text after A.

\roomfor{2}
Lead-in B, which says what the listing shows.\label{b-lead}
\testlisting{b}{30}
Text after B.

\roomfor{10}
Lead-in C.\label{c-lead}
\testlisting{c}{30}
Text after C.

\roomfor{29}
Lead-in D.\label{d-lead}
\testlisting{d}{30}
Text after D.

\roomfor{29}
Lead-in I, with line 28 empty.\label{i-lead}
\testblank{28}%
\testlisting{i}{30}
Text after I.

\roomfor{3}
Lead-in E, a file read by \verb|\inputminted|:\label{e-lead}
\testfile{e}{5}{|}{|}%
\inputminted[escapeinside=||]{text}{\jobname-e.txt}
Text after E.

\makeatletter
\ifdefined\nwbegincode
\makeatother
\roomfor{2}
Lead-in F, a code chunk:\label{f-lead}
\nwbegincode{1}\moddef{hello.py}\endmoddef\nwstartdeflinemarkup\nwenddeflinemarkup
\catcode37=14\relax\catcode126=13\relax\catcode36=3\relax\catcode38=4\relax\catcode35=6\relax\catcode94=7\relax\catcode95=8\relax\catcode13=5\relax\partopsep=0pt\relax\begin{minted}[xleftmargin=\codemargin,vspace=0pt,escapeinside=||]{python}
def main():|\label{fL1}|
    print("Hello, world!")

if __name__ == "__main__":
    main()|\label{fL5}|
\end{minted}
\nwendcode{}\nwbegindocs{2}\nwdocspar
Text after F.
\nwenddocs{}
\else
\makeatother
\fi

\roomfor{2}
Lead-in J, to code side by side:\label{j-lead}
\begin{didacticpair}{2}
  \rule{1pt}{8\baselineskip}\label{j-pair} & right column
\end{didacticpair}
Text after J.

\roomfor{9}
Lead-in L, to a short listing with room for eight of its ten lines.%
\label{l-lead}
\testlisting{l}{10}
Text after L.

\roomfor{6}
\section*{Heading K}\label{k-head}
First line of K\\
second line\\
third line\label{kL3}\\
fourth line.

\roomfor{2}
\begin{example}[Heading M]\label{m-head}
First line of M\label{mL1}\\
second line\\
third line\label{mL3}\\
fourth line.
\end{example}

\roomfor{6}
\section*{Heading N}\label{n-head}
Table N lists the rows.
\begin{longtable}{ll}
\caption{A caption of several lines, several lines, several lines, several
  lines, several lines, several lines, several lines, several
  lines.}\label{n-tab}\\
a & b\\ a & b\\ a & b\\ a & b\\
\end{longtable}

\roomfor{8}
\begin{enumerate}
\item Lead-in G, in an item:
  \testlisting{g}{5}
\end{enumerate}

\clearpage
\section*{Results}
\begin{itemize}
\item[] A short listing doesn't break and stays with its lead-in:
\testcheck{A first line with lead-in}{a-lead}{=}{aL1}
\testcheck{A first line on the next page}{a-lead}{=}{aL10}
\item[] A long listing keeps three lines at each end:
\testcheck{B lead-in with line 3}{b-lead}{=}{bL3}
\testcheck{C breaks}{cL1}{<}{cL30}
\testcheck{C lead-in with line 3}{c-lead}{=}{cL3}
\testcheck{C line 28 with the last}{cL28}{=}{cL30}
\testcheck{D breaks}{dL1}{<}{dL30}
\testcheck{D line 28 with the last}{dL28}{=}{dL30}
\testcheck{D line 27 with the first}{dL1}{=}{dL27}
\testcheck{I line 27 with the last}{iL27}{=}{iL30}
\testcheck{I breaks}{iL1}{<}{iL30}
\item[] minted, noweb and lists:
\testcheck{E lead-in with the listing}{e-lead}{=}{eL5}
\makeatletter
\ifdefined\nwbegincode
\makeatother
\testcheck{F lead-in with the header and code}{f-lead}{=}{fL1}
\testcheck{F the whole chunk}{f-lead}{=}{fL5}
\else
\makeatother
\fi
\testcheck{G first and last line together}{gL1}{=}{gL5}
\testcheck{J lead-in with the pair}{j-lead}{=}{j-pair}
\item[] Limits on the empty space and headings:
\testcheck{L breaks, rather than leave the room empty}{lL1}{<}{lL10}
\testcheck{L lead-in with line 3}{l-lead}{=}{lL3}
\testcheck{K heading with three lines}{k-head}{=}{kL3}
\testcheck{M block heading with its third line}{mL1}{=}{mL3}
\testcheck{N heading with its long table}{n-head}{=}{n-tab}
\end{itemize}

\end{document}
